% Generated by D:/tools/scratch_qdd/mmrec/render.py from data_g.py - edit the data, not this file.
\subsection{G. 연구 일지 --- 09-28 15시 \textasciitilde{} 09-30 21시 (KST, 날짜순)}\label{rec:G}
\noindent\textbf{한 줄:} 공개 데이터를 머리 시점·정확도 우선으로 모으고, 시뮬을 사실적 L8S로 다시 만들어 양산했다. 비율 실험(0.75 상한, 반복 2배)과 L8S 양 실험 뒤 확인판 E-C35가 관문을 못 넘었지만, 원인 진단 뒤 사용자 결정으로 본 35B(202,397행, 3에폭)를 09-30 15:57에 시작했다. 같은 날 다음 데이터 L9와 VLA = 좁은 조이스틱 결합 설계를 정했다.

이 절은 날짜순 일지다. 앞 A--F 줄기는 주제별이고, 여기서는 09-28 오후부터 09-30 밤까지 사용자가 말한 것, 정한 것, 일어난 일, 결과를 시각 순서대로 적는다(user-log 163--187, 189--218; 188은 C13). 시각은 KST이고, 사용자 말은 대부분 통제자(0300a0 세션)를 거쳐 전달되었다.

\paragraph{G1. 공개 데이터를 모두 모으되, 정확도를 떨어뜨릴 수 있는 것은 아예 쓰지 않는다}
{\small 09-28 (일) 15:3x -- 16:4x KST}
\begin{itemize}
  \item \textbf{사용자} (ul 163, 15:3x KST): \intent{고고 우리 이 데이터를 위한 할당량은 2TB로 일단 해줄게}
  \item \textbf{사용자} (ul 163, 15:3x KST): \intent{아니다 3TB로 해줄게 그렇다고 꽉차간다면서 너가 막 지우짐 ㅏㄹ고 나한테 물어봐야해 꼭}
  \item \textbf{사용자} (ul 163, 15:3x KST): \intent{전체 3TB로 해주고 하다가 2TB 넘어가면 나한테 일단 알려줘}
  \item \textbf{그래서(결정):} AgiBot World Beta를 점 라벨로 변환 시작(받기 $\rightarrow$ 추출 $\rightarrow$ 원본 삭제의 흘려 처리). AgiBot+RH20T 합산 3 TB, 2 TB를 넘으면 사용자에게 알림, 2.7 TB에서 내려받기 중지 후 문의. 공간 확보용 임의 삭제 금지. CC BY-NC-SA라 연구 한정, Astra 전송 금지.
  \item \textbf{사용자} (ul 164, 16:0x KST): \intent{지금 가용가능한 적합한놈을 다 써보자. 보든 처리 하기로한것도 다 거치는거 맞지?}
  \item \textbf{그래서(결정):} AgiBot·BEHAVIOR·RB2·RB3·MolmoBot·ManiSkill·CALVIN·RH20T·DROID를 모두 D 점 라벨로 변환 시도. 공통 처리 = 투영 $\leq$5 px(SAM 검증), 이름 검사, 점 라벨만(xyz·깊이·의도 없음), source·license 태그, 3인칭 태그.
  \item \textbf{사용자} (ul 165, 16:4x KST): \intent{그 정확도가 떨어질 수 있는 데이터는 아예 안쓰는게 맞아 알지?}
  \item \textbf{그래서(결정):} 원칙: 3인칭 RH20T·DROID, 관문 경계 출처(RB1), 5 px 초과·이름 검사 실패 행, PixMo 제외. 새 데이터는 사전등록한 8B 비교에서 모든 세트 비열등일 때만 채택.
\end{itemize}

\paragraph{G2. 공개 비율 비교는 L8 밖의 범용 평가로 --- 일반 평가 G}
{\small 09-28 (일) 17:1x KST}
\begin{itemize}
  \item \textbf{사용자} (ul 166, 17:1x KST): \intent{그냥 이건 생각이고 저 비교가 불필요한게 아닌데 만약에 비교를 할때는 l8기반의 환경이 아닌 다른 무언가여야한다는 얘기애}
  \item \textbf{그래서(결정):} 공개 데이터의 비율·효과는 L8-X만으로 판정하지 않는다. 일반 평가 G = 학습에서 뺀 공개 데이터 과제·기체의 물체 가리키기(AgiBot·BEHAVIOR·MolmoBot 보류), RB2 실물 보류, Where2Place·RefSpatial-Bench. 판정 = G 개선 + L8-X 비열등.
\end{itemize}

\paragraph{G3. 공개 데이터에 단계 의도를 붙여 사전학습 $\rightarrow$ 우리 로봇 포스트 트레인?}
{\small 09-28 (일) 20:4x KST}
\begin{itemize}
  \item \textbf{사용자} (ul 167, 20:4x KST): \intent{일단 저렇게 하고 포스트 트레인으로 우리 로봇에 맞게 하면 되지 않을까?}
  \item \textbf{해 봄:} E-STAGE8(8B, 사전등록): 1단계 넓은 학습(공개 데이터 + 기체 무관 단계 의도) $\rightarrow$ 2단계 L8-X 포스트 트레인 대 한 번에 섞기.
  \item \textbf{결과:} ul 186 무렵 등록 규칙대로 REJECT --- 두 단계 방식이 한 번에 섞기보다 낫지 않았다. 본 학습은 한 번에 섞는다.
\end{itemize}

\paragraph{G4. 공개 점 비율 재확인, MolmoBot 어안 제외, 3인칭은 평가로}
{\small 09-28 (일) 21:1x -- 22:5x KST}
\begin{itemize}
  \item \textbf{사용자} (ul 168, 21:1x KST (선택지 응답)): \intent{한 번 더 돌리기 (추천)}
  \item \textbf{그래서(결정):} E-OPRATIO8 첫 결과(p75 G 0.303 $\rightarrow$ 0.684, OOD-O·dev 개선, 표본 부족으로 비열등 미증명)에 p0·p75를 시드 1로 한 번 더. Where2Place 비열등 하한은 $-$0.03 $\rightarrow$ $-$0.10으로 완화(같은 자리 두 번째 선택).
  \item \textbf{사용자} (ul 169, 22:4x KST): \intent{광각 빼자 그럼 광각 아닌 예시 데이터 보여줘봐봐}
  \item \textbf{그래서(결정):} MolmoBot gopro 계열 광각(어안) 행 제외(Franka 검증 146, obj\_pixel 619; 필터이지 삭제 아님). 카메라별 view 태그 오표기(3인칭인데 head camera) 정정.
  \item \textbf{사용자} (ul 170, 22:5x KST): \intent{제대로 된 평가를 통해 제안의 123을 선정}
  \item \textbf{해 봄:} E-VIEW8(8B, 두 시드): A0 머리만 / A1 +3인칭 물체·놓을 곳 점 / A2 +3인칭 자기 손 점 / A3 +RH20T·DROID.
  \item \textbf{결과:} ul 186 무렵 A1·A2 REJECT $\rightarrow$ 머리 시점만. A3은 DROID 보류(G7)로 열지 않음.
\end{itemize}

\paragraph{G5. 약한 고리를 본 학습과 병행, RB2 보류분 일부를 학습으로}
{\small 09-29 (월) 00:1x -- 00:4x KST}
\begin{itemize}
  \item \textbf{사용자} (ul 171, 00:1x KST): \intent{지금 남는시간에 동시에 진행할 수 있는 약한고리 해결할 수 있는 방법 있으면 동시에 돌려봐도 좋음 최대한 대규모 학습 시간에는 영향이 없도록}
  \item \textbf{그래서(결정):} 본 학습에 즉시 양보하는 조건으로 병행: 한계 지도 교란 주입, 깊이 결손 대책, 폐루프 세트당 30편, RB2 실물 스테레오 깊이로 D 변환기 3D 오차, SAM 필터 쏠림 점검, 에폭 중간 체크포인트 평가.
  \item \textbf{사용자} (ul 172, 00:4x KST): \intent{1번은 바로 해도 문제는 없을 것 같음. 2번은 할 예정이라는 거지? 3번은 그냥 빼는게 맞고}
  \item \textbf{그래서(결정):} RB2 G 5,485행 중 약 1,500행만 G로 남기고 약 4,000행을 학습으로(에피소드 단위·해시 층화, 결과 보기 전). SAM 탈락·광각·이름 object·AgiBot 손끝·CALVIN은 제외 유지.
\end{itemize}

\paragraph{G6. 시뮬 품질: 도형 장면에서 사실적 장면 b4로}
{\small 09-29 (월) 01:3x -- 03:5x KST}
\begin{itemize}
  \item \textbf{사용자} (ul 173, 01:3x KST): \intent{퀄리티 올릴 준비 해놔보자}
  \item \textbf{사실 확인:} 집계: b3 대상의 81 \%가 원시 도형, 가구의 92 \%가 평판 탁자, 메시 가구·실물 물체·방·조명 네 요소를 모두 갖춘 편 0 \%.
  \item \textbf{사용자} (ul 174, 01:4x KST): \intent{단순한걸로 조합한 기존것듯 몇개는 있어도 되는대 비중이 많을 필요는 없음}
  \item \textbf{사용자} (ul 174, 01:4x KST): \intent{b4를 기다렸다 포함 (추천)}
  \item \textbf{사용자} (ul 175, 02:2x KST (음성 입력)): \intent{도용으로 만든 옛날 방식 같은 경우에는 15\% 미만 수준을 유지하게 해줘}
  \item \textbf{사용자} (ul 176, 02:3x KST (음성 입력)): \intent{같은 것도 좀더 늘릴 수 있으면 좋겠고 그리고 뭐 밀기나 아니면 은 사람 열기나 그 외에 다양한 테스트들을 할때 그런 테스트들에 적합한 기물들도 가져와 가지고 이미 움직임 설정이나 이런 거다돼 있는 기물들 가져와 가지고 하는 것도 좋을 것 같아}
  \item \textbf{사용자} (ul 177, 02:5x KST (선택지 응답)): \intent{75\% 채택 (추천)}
  \item \textbf{결과:} E-OPRATIO8 두 시드 합동: 범용 G 적중 +0.363 [+0.344, +0.383], L8-X 5세트 비열등·dev/OOD-H/OOD-O 유의 개선, OOD-T·OOD-H-lift는 꼬리 분산으로 비열등 미증명. 등록 판정 NOT\_QUALIFIED였으나 사용자 결정으로 본 학습 공개 점 75 \% 채택.
  \item \textbf{사용자} (ul 178, 03:1x KST (음성 입력)): \intent{조도 같은거 다르게 광원 위치나 이런것도 그렇고 다양하게 만들면 좋겠음}
  \item \textbf{사용자} (ul 179, 03:2x KST): \intent{아예 없애진 않아도 비중이 그리 많을 필요없겠는데}
  \item \textbf{사용자} (ul 179, 03:2x KST): \intent{한 10프로로 해두됨}
  \item \textbf{사용자} (ul 180, 03:3x KST): \intent{단순도 조명 배경 바뀌는 건 여전하게 하는 걸루}
  \item \textbf{사용자} (ul 181, 03:5x KST): \intent{그 캠을 건드리거나 할때 실제 urdf상에 있는 목을 돌리던지 하는 방식으로 해결해 보기를 바람 b4 는 1,2,4 적용해보면 좋을듯요}
  \item \textbf{사용자} (ul 181, 03:5x KST): \intent{근데 목 무작위화도 한 15프로 정도로만}
  \item \textbf{그래서(결정):} 사실적 묶음 b4: 실물 물체와 어지러운 물체를 가구 면에, 방 배경 전 편, CC0 PBR 재질·HDR, 조명(밝기·위치·개수·색) 무작위화, 센서 현실감(노출·잡음·흐림·JPEG). 도형 장면 <15 \%(단순 묶음도 b4 조명·배경으로 재생성), 쟁반 위 놓기 $\leq$10 \%, 도착지는 바구니·통·상자·그릇·다른 가구 면·진열 빈자리·물체 옆·상품 위로. 묶음 = 단순 / 사실적(b4) / 서랍(b3d). 시점 변화는 URDF 목 관절(pan·tilt)로, 편의 약 15 \%만.
\end{itemize}

\paragraph{G7. 기물 상호작용 30종 이상, 새 이름 L8S, DROID는 마지막}
{\small 09-29 (월) 04:1x -- 04:5x KST}
\begin{itemize}
  \item \textbf{사용자} (ul 182, 04:1x KST): \intent{바구니도 좋고 책장도 좋고 테스크를 수행할때 다양한 기물과 상호 관계가 많이 있음 좋겠어 상자에 넣거나 냉장고를 열거나 이런거 전부 다 말이야.}
  \item \textbf{사용자} (ul 183, 04:2x KST): \intent{이것보다도 더 많아야해}
  \item \textbf{그래서(결정):} 기존 집기·놓기 어휘로 되는 기물 상호작용 과제 30종 이상(용기·싱크·장바구니·열린 수납 안에 넣기, 선반 칸 놓기·꺼내기, 관계 놓기, 속성 지정 집기, 모으기·분류·되돌리기 등). 관절 기물 다단계 과제는 새 어휘가 필요해 다음 판.
  \item \textbf{사용자} (ul 184, 04:3x KST): \intent{오케이 좋은 데이터가 l8을 통해만들어질 수 있게 된거지? 이버전 l8을 l8s 라하자}
  \item \textbf{사용자} (ul 184, 04:3x KST): \intent{목 무작위화는 좌우 상하에 영향을 미치는 걸로 합시다 목 롤 은 뭐 딱히 건들지 말구}
  \item \textbf{사용자} (ul 184, 04:3x KST): \intent{상하 10도 좌우 15도 단 작업을 할때 목표와 상품이 보이는 상황이 되는경우에만 해당함}
  \item \textbf{그래서(결정):} 새 판 이름 L8S = L8S-real(b4) + L8S-simple(<15 \%) + L8S-drawer(b3d). 목은 pan $\pm$15$^\circ$·tilt $\pm$10$^\circ$, 롤 고정, 대상·도착지가 화면에 보일 때만. 옛 b1--b3는 L8-X로 보관.
  \item \textbf{사용자} (ul 185, 04:5x KST): \intent{일단 드로이드는 제일 마지막으로 미루고 다른거 다하고 방법을 찾아보고 안되면 유기하는걸루}
  \item \textbf{그래서(결정):} DROID 3인칭 물체 점 변환(수율 29/1,980편) 정지·상태 저장. 결국 본 학습에 들어가지 않음.
\end{itemize}

\paragraph{G8. 비교 판정 규칙 개정 --- 시드 흔들림으로 여유 보정, 중앙값·실패율로}
{\small 09-29 (월) 05:4x KST}
\begin{itemize}
  \item \textbf{사실 확인:} 같은 데이터·다른 시드끼리도 L8-X 평균 상한 +2 mm 비열등을 35번 중 한 번도 통과하지 못했다 --- 규칙이 시드 잡음보다 좁았다.
  \item \textbf{사용자} (ul 186, 05:4x KST (선택지 응답)): \intent{여유 보정 + 강한 지표 (추천)}
  \item \textbf{그래서(결정):} ni\_judge: 비열등 여유 = 같은 모델 시드 쌍(A/A)의 흔들림 폭, 평균 대신 중앙값과 20 mm 초과 실패율. E-VIEW8(머리 시점만)·E-STAGE8(한 번에 섞기)은 등록 규칙대로 확정.
\end{itemize}

\paragraph{G9. 실패 회복 데이터는 이번에 넣지 않는다 --- 뒤로 미룬 문제 목록}
{\small 09-29 (월) 오후 KST}
\begin{itemize}
  \item \textbf{사용자} (ul 187): \intent{일단 넣지 않는걸로 하고 그거 하나 만들어 놓을까 나중에 해결해야하는 문제들 이라는 이름으로}
  \item \textbf{그래서(결정):} 회복 데이터(실패 인지·재시도)는 본 35B에서 뺀다(역효과 위험). 넣으려면 8B 비교에서 전 세트 비열등이어야 한다. \texttt{docs/\allowbreak{}later\_\allowbreak{}problems.\allowbreak{}md}에 뒤로 미룬 문제를 한곳에 모은다(지금 22항목).
  \item \textbf{참고:} 같은 날 18:5x KST의 ul 188(VLA = 늘 쓰는 조이스틱)은 C13에 있다. 그 설계는 09-30 밤 G21에서 좁은 조이스틱으로 다시 정해졌다.
\end{itemize}

\paragraph{G10. L8S 시범 재검수 기준: 수율 70 \% $\rightarrow$ 50 \%, 왼손목 튐 병행}
{\small 09-29 (월) 20:4x -- 22:0x KST}
\begin{itemize}
  \item \textbf{사용자} (ul 189, 20:4x KST): \intent{수율을 70프로이상이면 통과하는걸루}
  \item \textbf{사용자} (ul 190, 20:4x KST): \intent{빈카드로 왼손 목 튀는 문제도 해결을 해봐야할듯}
  \item \textbf{그래서(결정):} 놓는 순간 left joint7이 0.073/0.081 rad 튀는 문제(손목이 토크 한계 5.1 Nm에 걸렸다가 한꺼번에 풀림)를 x2 빈 카드에서 진단. 팔 걸음 $\leq$0.04 rad 규칙은 그대로.
  \item \textbf{사용자} (ul 191, 22:0x KST): \intent{수율 50프로면됨}
  \item \textbf{결과:} 시범6 공식 수율 50 \%(20/40) --- 통과. L8S 양산 22:50 KST 시작(계획 8,917편).
\end{itemize}

\paragraph{G11. 본 학습 전 확인판 E-C35, 공개 비율 0.75 + 반복 상한 2배, L8S 전량}
{\small 09-29 (월) 22:1x -- 23:4x KST}
\begin{itemize}
  \item \textbf{사용자} (ul 192, 22:1x KST): \intent{L8S 양산 초반 일부로 본 학습 전에 소규모 확인 판을 한 번 돌리자}
  \item \textbf{해 봄:} E-C35 사전등록 cac60bd: L8S 양산 앞 1,000편 + 공개 점으로 35B 대 f35\_d, 평가 L8-X dev·OOD-H·OOD-O + 새 물체 OOD-O58 + G. 이후 변경 1--4로 비율·두 팔·층화 방식을 고정.
  \item \textbf{사용자} (ul 193, 22:3x KST): \intent{1번먼저검사하고 사시 0으로 일단 하고 그다음에 지금 비중보다 더 늘리는게 줄이는게 좋은지도 검사하면 좋을듯 10\% +-로}
  \item \textbf{결과:} E-OPR2 1단계(학습 없음, 새 ni\_judge로 재판정, a35da09): 당시 recipe(공개/기본 0.59) 유지 --- p75 두 시드 대 p0이 L8-X 7세트·G 모두 비열등, p25·p50은 L8-X 실패율이 p75보다 나쁨.
  \item \textbf{사용자} (ul 194, 22:3x KST): \intent{고}
  \item \textbf{사용자} (ul 195, 23:0x KST): \intent{음 근데 그러면 우리도 75\%로 하는게 일단 좋지 않을까 싶은데? 나중에 해야할 투두에다가 저 최적의 오픈, 만든것 비율 알아봐야한다고좀 해줘봐봐}
  \item \textbf{사용자} (ul 196, 23:0x KST): \intent{지금 파이프라인에서 빠지는 것들은 쫙 검토해서 쳐내줘 예전것들이 남아있는게 좀많아}
  \item \textbf{결과:} 논문 6eaa0b0: Astra 증류 서술 $\rightarrow$ 시뮬 참값 + 공개 점 라벨, R/H·C$'$·M4 규칙·옛 마일스톤을 걷어내고 L8S 양산 $\rightarrow$ E-C35 $\rightarrow$ 본 35B $\rightarrow$ 평가 흐름으로.
  \item \textbf{사용자} (ul 197, 23:1x KST): \intent{그럼 딱 그만큼 맞춰서 고하자}
  \item \textbf{그래서(결정):} 공개/시뮬 = min(0.75, 2.0 $\times$ 공개 풀 / 시뮬 기본 행), 반복 상한 2.0배(ul 195의 0.75 고정을 대체). E-C35를 반복 효과 짝 비교로 확장: 팔 b(반복 없음, 시드 0·1) 대 팔 a(풀 부분집합을 반복해 0.75).
  \item \textbf{사용자} (ul 198, 23:2x KST): \intent{뭐 너가 올바른거 근거기반으로 정해줘봐}
  \item \textbf{그래서(결정):} 위임받은 결정: 본 35B는 L8S 계획 전량 8,917편(B안) --- 1,200편 안은 E-C35와 거의 같아 정보가 적고, 일반화에는 방·가구·물체·과제 조합의 다양성이 필요. 논문 평가 프로토콜(Inspect Robots + RD)은 유지.
\end{itemize}

\paragraph{G12. 765a 카드 두 장을 양산에, 13:20 반납}
{\small 09-30 (화) 01:0x -- 01:3x KST}
\begin{itemize}
  \item \textbf{사용자} (ul 199, 01:0x KST): \intent{765a에서 렌더가 된다고 확인된 카드는 GPU1·GPU3 이거 두장 잡아서 사용해볼래?}
  \item \textbf{사용자} (ul 199, 01:0x KST): \intent{빨리 확인해보자}
  \item \textbf{결과:} 카드별 렌더 시험: 두 장만 렌더 가능, 세 장은 Vulkan DEVICE\_LOST(학습은 가능). 5장 파드를 1장 파드 5개로 다시 잡아 렌더 되는 두 장(g1·g5)에 L8S 레인 추가.
  \item \textbf{사용자} (ul 200, 01:3x KST): \intent{저 2장은 1시 20분에 반납하는걸루}
  \item \textbf{그래서(결정):} 13:15까지 레인 정리, 13:20 파드 안 예약으로 자기 종료.
\end{itemize}

\paragraph{G13. 비율 확장 탐색 대 L8S 비중 실험 --- 순서를 두 번 바꿈}
{\small 09-30 (화) 01:2x -- 02:1x KST}
\begin{itemize}
  \item \textbf{사용자} (ul 201, 01:2x KST): \intent{고}
  \item \textbf{그래서(결정):} E-OPR3 사전등록 2a4f49a: 공개/시뮬 1.0·1.5·3.0(8B 시드 0, 공개 행은 1.5배 반복 고정, 시뮬 행을 줄여 비율을 맞춤).
  \item \textbf{사용자} (ul 202, 01:3x KST): \intent{아니 근데 시뮬을 더 적게쓰면 어떨까라기보다는 l8s의 비중이 과연영향을 미치는가가 궁금한데 그건 이미 근거가 있나?}
  \item \textbf{사용자} (ul 202, 01:3x KST): \intent{지금 78dc를 그럼 l8s의 비중으로 돌려 시뮬의 비중보다}
  \item \textbf{그래서(결정):} E-L8SW 사전등록 06accad: 공개 풀 32,365행 고정(반복 없음), 옛 L8-X 없이 L8S만 25·50·100 \%(C35 층화 1,000편에서 추출), 8B 시드 0. 추가 발언 ``시뮬비중은 그다음으로'' $\rightarrow$ OPR3는 보류가 아니라 뒤에 재개.
  \item \textbf{사용자} (ul 211, 02:0x KST): \intent{그럼 빨리 비율탐색 먼저하는걸루 그다음에 이어서}
  \item \textbf{결과:} E-OPR2 2단계($\pm$10 \%p, 시드 0): 0.49·0.59·0.69 사실상 같음 $\rightarrow$ 0.59 유지(0.49는 OOD-T 실패율만 조금 나쁨).
  \item \textbf{결과:} E-OPR3: 1.0·1.5·3.0으로 공개 비중을 올리면 L8-X가 나빠지고 G는 나아지지 않음 $\rightarrow$ 후보 없음.
  \item \textbf{사실 확인:} E-L8SW 도중 결함: 섞기 스크립트가 편 이름을 잘못 맞춰 첫 세 판이 L8S 0행(공개만)이 됨 $\rightarrow$ l25·l50 무효 처리·격리(변경 3), 다시 섞어 L8S 4,254 / 8,558 / 17,094행으로 재학습.
  \item \textbf{결과:} E-L8SW: L8S가 많을수록 새 물체 중앙 30.0 $\rightarrow$ 10.2 mm, L8S 보류 8.2 $\rightarrow$ 5.2 mm, OOD-H 실패 29 $\rightarrow$ 24 \%. 50 $\rightarrow$ 100 \%는 이득이 작음(포화). 엄격 규칙은 OOD-T·G Franka 여유 때문에 \texttt{해로움} 판정 --- 시드 하나라 방향 신호로만 읽음.
\end{itemize}

\paragraph{G14. 나중 할 일 목록에 올린 것들}
{\small 09-30 (화) 01:4x -- 02:0x KST}
\begin{itemize}
  \item \textbf{사용자} (ul 203, 01:4x KST): \intent{투두에다가 손목캠을 활용방안 찾기 어찌하면 성능 높아질지 비교 분석 이것도 적어놔봐봐}
  \item \textbf{사용자} (ul 204, 01:5x KST): \intent{투두에 오른손 쓰는것도 넣어둬 다음에 하는걸루}
  \item \textbf{사용자} (ul 205, 02:0x KST): \intent{일단 왼손은 쓰지 않고 원래 하던 방향대로 가는걸로 할게}
  \item \textbf{사용자} (ul 206): \intent{그런것도 다해서 투두에 적어놔보자 걱정되는점}
  \item \textbf{사용자} (ul 207): \intent{실측 데이터 확보, 회복능력, 아스트라와 비교, 깊이 노이즈 있ㄴㄴ거 어떻게 해서 리얼에서 완벽한 방안으로 노이즈 최저의 값을 얻을건지 ai워커에서 말이야. 이거 들 투두에 적어놔줘}
  \item \textbf{사용자} (ul 209): \intent{노노 그 투두에 그냥 로로로 충분한지 알아보기 만 적자}
  \item \textbf{그래서(결정):} later\_problems 15(손목캠 활용), 16(왼팔 --- 사용자 결정: 하지 않음), 17(손목캠 걱정점: 손목 영상 이상 강건성 평가를 본 35B 평가에 넣음), 1·9 보강(회복·실측), 18(Astra 비교, 유료 허락 뒤), 19(AI Worker 실물 깊이 잡음 최소화), 20(LoRA로 충분한지).
  \item \textbf{참고:} ul 206--210의 원래 시각(02:1x·02:3x)은 짐작으로 적은 오기였고, 실제는 01:4x--02:0x KST다(8e25f8c 정정). 이후 시각은 date로 찍는다.
\end{itemize}

\paragraph{G15. 본 35B는 3에폭, 체크포인트 1·1.5·2·2.5·3 / 토큰 절약 규칙}
{\small 09-30 (화) 01:4x -- 02:0x KST}
\begin{itemize}
  \item \textbf{사용자} (ul 208): \intent{중간 체크포인트 저장 좋아 3에폭까지 늘리는 것도 좋아 저장은 1, 1.5, 2, 2.5, 3 저장하는 걸루 하자. 평가는 일단 자동으로하는거는 2 -> 1.5 -> 2.5 순서루}
  \item \textbf{그래서(결정):} main35\_recipe: 3에폭, 다섯 곳 저장, 자동 평가 2 $\rightarrow$ 1.5 $\rightarrow$ 2.5(이후 1 $\rightarrow$ 3). 최적 에폭은 학습에 안 쓴 검증 분할로, 결과 전에 사전등록.
  \item \textbf{사용자} (ul 210): \intent{클로드 토큰 사용량 최적화는 되어있나? 각각의 에이전트의 성능을 떨어뜨리는 건 목적이 아니구}
  \item \textbf{그래서(결정):} 세션 간 중계는 핵심만, 받은 쪽은 events.log 한 줄. 폴링 금지, 로그는 필요한 줄만, 보고 8줄 이하, 긴 세션은 압축·인계. 에이전트 성능·추론 강도·모델 등급은 낮추지 않는다.
\end{itemize}

\paragraph{G16. 서랍을 넣었다가 뺐다, 브랜치 정리, 확인판 뒤 본 학습까지 자동으로}
{\small 09-30 (화) 03:3x -- 04:3x KST}
\begin{itemize}
  \item \textbf{사용자} (ul 212, 03:4x KST): \intent{서랍 넣지 문제 없는지 보구}
  \item \textbf{사실 확인:} 점검: 서랍 b3d 148편은 관절 기록(joints.npz)이 없어 팔 튐을 오프라인으로 잴 수 없고, 팔 관성 수정 전 코드·L8-X standard 렌더로 만들어졌다.
  \item \textbf{사용자} (ul 213, 04시경 KST): \intent{일단 서랍은 빼자 그럼}
  \item \textbf{그래서(결정):} 서랍은 본 35B와 E-C35 모두에서 뺌(E-C35 변경 8, 학습 전). 재생성 조건은 later\_problems 4·21.
  \item \textbf{사용자} (ul 214, 03:3x KST): \intent{브랜치는 올바른 데브 하나만 유지}
  \item \textbf{사용자} (ul 215, 04시경 KST): \intent{그냥 원래 계획대로 하자 바꾸지 말구}
  \item \textbf{사용자} (ul 216, 04:1x KST): \intent{E-C35 통과 $\rightarrow$ 최종 빌드 $\rightarrow$ 본 학습 기동}
  \item \textbf{사용자} (ul 216, 04:1x KST): \intent{잇도 자동 골려있게 해야하는거 알지?}
  \item \textbf{결과:} 본 35B 자동 체인 52c122b: C35 판정 NONINFERIOR $\rightarrow$ 양산 끝(PROD\_DONE) $\rightarrow$ 남은 편 빌드 $\rightarrow$ 섞기 $\rightarrow$ 78dc 4장 3에폭, 아니면 ALERT\_GATE\_FAIL로 멈춤. 드라이런 두 분기 정상.
  \item \textbf{사실 확인:} 806f3a2: 선행 빌드 1,590편 실측으로 L8S 한 편 = 약 17.9행(조종 11.0 + 보조 6.9) --- 이전 8.7행 추정을 정정. 그래서 공개 점은 반복 2배 상한에 걸림.
\end{itemize}

\paragraph{G17. E-C35 관문 FAIL $\rightarrow$ 원인 진단 $\rightarrow$ 소규모 재실행 없이 본 학습}
{\small 09-30 (화) 13시 -- 15:1x KST}
\begin{itemize}
  \item \textbf{결과:} E-C35 팔 b 대 f35\_d WORSE: L8-X dev 실패율 16.8 $\rightarrow$ 27.5 \%, OOD-H 17.3 $\rightarrow$ 22.8 \%, G MolmoBot Franka 0.87 $\rightarrow$ 0.45, RBY1 0.94 $\rightarrow$ 0.78, G 전체 0.74 $\rightarrow$ 0.66. 좋아진 것: 새 물체 58편 중앙 40.9 $\rightarrow$ 9.6 mm·실패 60 $\rightarrow$ 36 \%, BEHAVIOR 0.71 $\rightarrow$ 0.77, RB2 0.92 $\rightarrow$ 0.95, RefSpatial·Where2Place. 팔 a(반복) 대 b는 OOD-H만 조금 나쁨.
  \item \textbf{사실 확인:} 진단(재학습 없음): (1) G Franka 보류 300행은 전부 3인칭·어안인데, f35\_d는 이런 행 8,469개를 학습했고 지금 풀은 머리 시점 규칙(ul 165·169)으로 0행 --- 규칙과 안 맞는 세트. (2) RBY1은 C35에서 280행뿐(f35\_d 4,814행). (3) L8-X 차이는 도메인과 양 --- f35\_d는 L8-X 조종 47,724행, C35는 L8S 조종 10,813행(4.4배 적음).
  \item \textbf{사용자} (ul 217, 13\textasciitilde{}14시 KST): \intent{그냥 본 학습에서 일단 수정을 먼저 어떻게 할 건지 계획 세우고 수정해서 본 학습 그냥 바로 들어가자 소규모 더 돌리지 말고}
  \item \textbf{사용자} (ul 217, 13\textasciitilde{}14시 KST): \intent{본 학습 바로 시작을 하자}
  \item \textbf{사용자} (ul 217, 13\textasciitilde{}14시 KST): \intent{수정안이 안됐으면은 시작은 안하는데 수정안 다 되면은 시작을 하자}
  \item \textbf{그래서(결정):} 수정안: (a) G Franka는 판정 제외·보고만 (b) RBY1 검증 통과분 추가 --- 확인해 보니 700행이 이미 풀에 전부, 추가분 없음 (c) 옛 L8-X는 넣지 않음 (d) 게이트 재판정 없이 기동. prereg\_main35 변경 1(관문 이탈 기록).
  \item \textbf{사용자} (ul 218, 15시경 KST): \intent{Franka를 판정에서 빼자 수동함}
  \item \textbf{사용자} (ul 218, 15시경 KST): \intent{내가 못한다니까 터미널에서 허락을 못 눌러}
  \item \textbf{결과:} d790efe: 체인 GATE\_OVERRIDE, 판정 스크립트 주 결론에서 Franka 제외(with\_franka 병기). 고리 V 설명 파서 버그(strip이 고리 설명을 모름) 고침 2fbdc54·eb9df85 $\rightarrow$ 고리 1,216행 변환.
\end{itemize}

\paragraph{G18. 본 35B 기동 --- 첫 판을 멈추고 기록을 늘려 처음부터 다시}
{\small 09-30 (화) 15:2x -- 16:0x KST}
\begin{itemize}
  \item \textbf{결과:} 15:27 KST 첫 기동(검증 분할 빌드가 같은 고리 버그로 한 번 실패 뒤 재기동). BUILD\_DONE 202,397행 = L8S 139,603(고리 V 1,216 포함, 서랍 없음, 옛 L8-X 없음) + 공개 점 62,794(채택 풀 31,397 $\times$ 반복 2.0, 공개/기본 0.45, 시뮬 약 69 \%). 검증 L8S 1,764행.
  \item \textbf{사용자 말(user-log 번호 없음)} (09-30 15:5x KST (통제자 경유, events.log)): 학습 관련 모든 정보를 다 담아야 한다, 체크포인트는 촘촘히.
  \item \textbf{그래서(결정):} 첫 판(약 30분)을 멈추고 보존, 7d91cd6으로 15:57 KST 처음부터 재기동. 추가 기록: run\_meta(데이터 sha256, 묶음·출처별 행, 판본, GPU, 인자), 기울기 노름(평균·최대), 묶음별 손실(시뮬 조종·시뮬 보조·공개), 어댑터 0.25에폭마다(12개). 평가 이름(1·1.5·2·2.5·3)과 순서 2 $\rightarrow$ 1.5 $\rightarrow$ 2.5 $\rightarrow$ 1 $\rightarrow$ 3 그대로, 0.25 단위는 별도 대기열.
  \item \textbf{진행 중·대기:} \tentative{3에폭 약 25,300걸음, 78dc H200 4장, 끝 약 10-01 13--14시 KST. 최적 에폭 = L8S 보류 검증 20 mm 초과 실패율 최저.}
\end{itemize}

\paragraph{G19. 본 35B 초기 체크포인트 --- 옛 도메인 차이가 줄어든다}
{\small 09-30 (화) 저녁 KST}
\begin{itemize}
  \item \textbf{진행 중·대기:} \tentative{ep0.25: 새 물체 중앙 9.5 mm·실패 35.5 \%, L8-X dev 실패 23.7 \%, OOD-H 25.0 \%.}
  \item \textbf{진행 중·대기:} \tentative{ep0.5: 새 물체 9.5 mm·34.8 \%, dev 21.4 \%, OOD-H 18.6 \%, G 0.67(Franka 포함). f35\_d는 dev 16.8 \%·OOD-H 17.3 \%.}
  \item \textbf{참고:} 학습이 진행될수록 L8-X 차이가 줄어드는 것은 C35 손해가 도메인보다 양 때문이라는 해석을 지지한다(아직 초기값).
\end{itemize}

\paragraph{G20. 다음 데이터 L9 --- L8S의 5배, 모든 환경이 서로 달라 보이게}
{\small 09-30 (화) 오후 -- 20:3x KST}
\begin{itemize}
  \item \textbf{사용자 말(user-log 번호 없음)} (09-30 오후 KST (0300a0 세션, 설계 문서 머리말)): l8s와 완전히 다른 느낌, 같은 데이터셋에서 온 게 맞나 할 정도로 차이. 선반장·식탁·낮은 책상 등 다른 환경, 우산 꽂기·컵 정리·정렬·접시 쌓기 등 새 과제, 목 각도·조명 다양화, 왼팔 사용, 양은 L8S의 5배.
  \item \textbf{사용자 말(user-log 번호 없음)} (09-30 오후 KST (설계 문서 §6)): 모든 환경은 l9의 환경은 진짜 각기 다른 환경으로 보여졌으면 좋겠어.
  \item \textbf{그래서(결정):} 설계 확정(spec 2026-09-30-l9-diverse-datagen-design, 계획 14fc475): 성공 편 $\geq$38,000, 과제 $\geq$75, 물체 $\geq$6,000, 가구 $\geq$1,400, 방 $\geq$190, 재질 $\geq$650, HDRI $\geq$140, 조명 계열 $\geq$10, 목 자세 100 \% 편에서 변화. 환경 8계열(선반장·식탁·낮은 탁자·주방·현관·사무실 책상·가게 진열대·작업대). 왼팔·오른팔 둘 다(출력에 hand 필드). 환경 조합 해시 중복 0 + 머리 f0 SigLIP 거의 같은 장면 비율이 L8S보다 낮을 것.
  \item \textbf{그래서(결정):} 움직임은 사람처럼: 두 구간 최소 저크 뻗기, 물체 크기에 맞춘 벌림, 자세 선호, 편마다 다른 조작자 스타일. 파지는 맞잡기 표본기 + 시뮬 파지 시험, 충돌 회피는 cuRobo V2(Apache)만. 비현실 모드는 만들되 later\_problems 22 비교 전까지 섞지 않음. 2차로 두 팔 협업 15--20 \%. 다른 공개 로봇은 로봇에 달린 머리 카메라가 있을 때만(3인칭뿐인 Franka·UR5e 등 제외).
  \item \textbf{결과:} 에셋 1차(18:42--20:31): 재질 1,307(CC0) + HDRI 501, 가구 862(THOR·Objaverse, 실제 높이로 재조정), 물체 1차 1,917 정착 통과, ProcTHOR 방 289, 용기 1차 관문 210 + L8S 35 = 245. 구현 진행 중.
\end{itemize}

\paragraph{G21. VLA 결합 설계 브레인스토밍 --- 명령 권한은 상위, VLA는 좁은 조이스틱}
{\small 09-30 (화) 저녁 KST (사양 문서 전)}
\begin{itemize}
  \item \textbf{사용자 말(user-log 번호 없음)} (09-30 KST (통제자 경유, board NOW.md 1-1 원문 요지)): 기존 VLA는 안의 VLM(3\textasciitilde{}8B)이 계획·판단·움직임·명령을 전부 떠안았다. 그 크기의 VLM에게는 너무 무거운 짐이고, VLM에서 action expert로 넘어가며 정보도 손실된다. 그래서 계획·판단·명령은 전부 상위 LLM이 맡고, VLA에게는 조이스틱이라는 아주 단순한 일(상위 명령 기준의 미세 조종)만 준다.
  \item \textbf{그래서(결정):} 구조의 핵심 가설 H-J(논문 서론·방법에 반영): 짐을 덜어 준 VLA가 같은 크기의 단독 VLA보다 낫다. VLA 설계·학습·평가는 VLA가 스스로 계획하게 만드는 쪽으로 가지 않는다 --- 입력 목표는 상위 명령, 출력은 허용 범위 안의 보정 + 이상 신호뿐. 결합 실험은 짐을 덜어 준 VLA 대 단독 VLA를 직접 비교하도록 짠다.
  \item \textbf{그래서(결정):} 권한: 상위가 명령 권한을 모두 가진다. VLA = 좁은 조이스틱 --- 허용 범위 안의 보정, 쥐는 시점, 작은 회피, 부딪힌 뒤 재정렬 + 이상 신호. (C13의 \texttt{늘 섞기} 틀을 대체)
  \item \textbf{그래서(결정):} 지연: 사건 기반 조기 재질의 + RTC식 겹침. 상위는 도착 전에 다음 명령을 미리 내도록 학습 --- 행 = 도착 전 $\Delta$ 0--3 s 프레임, 실행 중인 명령, FK 도착 표시를 입력으로 그림 $\rightarrow$ 다음 명령 + \texttt{지금 미리 내도 안전} 표지. VLA가 실제로 도착할 지점의 자기 예측은 작은 보조 후학습 과제로만, 제어 출력에는 넣지 않음.
  \item \textbf{그래서(결정):} VLA 학습 단계: 혼자(교사 보정) $\rightarrow$ 측정한 상위 오차·지연·도중 교체 주입 $\rightarrow$ 실제 상위 명령 위 $\rightarrow$ 상위와 번갈아. 명령 형식과 픽셀$\rightarrow$xyz 코드는 공유.
  \item \textbf{그래서(결정):} 채택 기준(사용자 결정): 결합 손실(두 오라클 조합 상한 중 작은 값 $-$ 결합 성공률) $\leq$ 5 \%p, 그리고 인터페이스 탓 실패 $\leq$ 전체 실패의 10 \%.
  \item \textbf{그래서(결정):} 최종 모델: 다음 본 학습에서 기본 모델부터 L8S + L9 + 선발행 행을 섞어 학습. 지금 어댑터에 이어 학습은 빠른 가능성 확인으로만.
\end{itemize}

\paragraph{G22. GPU 배분과 폐루프 확인 E-M35CL 시작}
{\small 09-30 (화) 20:3x -- 21:2x KST}
\begin{itemize}
  \item \textbf{그래서(결정):} 사용자 GPU 배분(20:47): 학습 = 78dc 0--3, 평가 = 7a2a GPU2 한 장(판정 대기열 2·1.5·2.5·1·3 + 0.25 대기열), VLA 개발 예약 = x2 GPU0·x3 GPU0(비어 있으면 폐루프 서버가 빌려 쓰고 요청 파일이 생기면 비움), L9 양산 = 7a2a GPU0·1·3 + x2 GPU1(양보 파일).
  \item \textbf{해 봄:} E-M35CL 사전등록 d8ba238(보고 전용): OOD-O58 + L8S 보류 약 160편 $\times$ \{기본, 어두운 따뜻한 조명, 머리 +10$^\circ$/+15$^\circ$\}, f35\_d $\rightarrow$ 0.5 $\rightarrow$ 1 $\rightarrow$ \ldots{} $\rightarrow$ 3, 모든 편 머리·양손목 영상. 변경 1: 빌린 서버 카드는 요청 시 즉시 비우고, 서버 오류 편은 실패로 세지 않고 다시 돌림.
  \item \textbf{진행 중·대기:} \tentative{21:19 KST: 레인 9개(7a2a GPU0·1·3), f35\_d 82편 끝·오류 0, ep0.5 병합 준비. 렌더 사용률 14--20 \%(CPU 병목).}
  \item \textbf{결과:} 논문 갱신(이 판): 본문은 지금 방법과 살아 있는 근거만 남기고(옛 행 수·비율·늘 섞기·끝난 계획 줄 삭제), 이력은 이 일지로.
\end{itemize}
